\section{Trabalhos Relacionados}

\subsection{Bases e tarefas em citologia cervical}

Bases como Herlev e SIPaKMeD são frequentes em estudos de classificação
cervical \cite{jiang2023systematic,fang2024systematic}. A SIPaKMeD reúne
células extraídas de esfregaços de Papanicolau para classificação por atributos
e imagens \cite{plissiti2018sipakmed}. Embora úteis para comparação, bases de
células isoladas não reproduzem integralmente sobreposição, inflamação,
variação de fundo e outros fatores encontrados em exames convencionais.

A CRIC Cervix foi construída a partir de 400 imagens de citologia convencional
e contém 11.534 células classificadas por citopatologistas
\cite{rezende2021cric}. A coleção inclui NILM, ASC-US, LSIL, ASC-H, HSIL e
SCC, preservando desafios de lâminas reais, como células sobrepostas e
elementos inflamatórios. Por isso, é útil para estudar variação entre imagens,
desde que a dependência entre células da mesma imagem seja considerada.

\subsection{Aprendizado profundo para classificação cervical}

Estudos recentes reportam desempenho elevado com redes profundas em Herlev,
SIPaKMeD e outras bases de citologia cervical. Kaur et al. compararam modelos
de transferência e reportaram 95\% de acurácia com ResNet50 em Herlev binário
\cite{kaur2025comparison}; Mehedi et al. obtiveram 98,53\% de acurácia com
CCanNet em SIPaKMeD de cinco classes \cite{mehedi2024lightweight}; e Lima et
al. reportaram 98,35\% em Herlev binário e 99,73\% em SIPaKMeD binário com
\emph{ensembles} \cite{lima2025ensemble}. Esses resultados, embora fortes, não
são diretamente comparáveis ao presente estudo por diferenças de base,
preparação, número de classes, métrica principal e unidade de particionamento.

Em citologia convencional, Rodríguez et al. reportaram sensibilidade 0,688 e
especificidade 0,762 em esfregaços de Papanicolau, indicando maior dificuldade
que em bases de células isoladas ou citologia em meio líquido
\cite{rodriguez2024automated}. Este trabalho se diferencia pelo protocolo:
CRIC, separação por imagem, validação cruzada agrupada, limiares definidos na
validação, incerteza por bootstrap agrupado, comparação local com ResNet50 e
análise de erro por nomenclatura Bethesda.
